% !TeX root = main.tex
\usepackage{amsmath}
\usepackage{ascmac}
\usepackage{amsthm}
\usepackage{amsfonts}
\usepackage{latexsym}
\usepackage{bm}
\usepackage{mathtools}
\usepackage{empheq}
\usepackage{amssymb}
\usepackage[dvipdfmx]{graphicx}
\usepackage{tikz}
\usepackage{textcomp}
\usepackage{enumitem}
\usepackage[most]{tcolorbox}
\usepackage[dvipdfmx]{xcolor}
\usepackage{environ}
\usetikzlibrary{intersections,calc,arrows.meta}

% ========================================
% 解答表示の切り替え（27行目だけ書き換える）
% ========================================
%   \showanswerfalse  … 生徒用（空欄・解答非表示）
%   \showanswertrue   … 先生用（青字で解答表示）
% 切替後: ターミナルで make all  または  Cmd+S で保存
\newif\ifshowanswer
\showanswertrue

% 行間を少し広げる
\linespread{1.15}

% 分数を常に大きく表示
\everymath{\displaystyle}

% 数式環境前後の空白を調整（ドキュメント開始時に適用）
\AtBeginDocument{%
  \setlength{\abovedisplayskip}{1pt}%
  \setlength{\belowdisplayskip}{6pt}%
  \setlength{\abovedisplayshortskip}{1pt}%
  \setlength{\belowdisplayshortskip}{3pt}%
}

% 平均記号のバーを長くするコマンド
\newcommand{\mean}[1]{\overline{#1}}

% 分数の穴埋めコマンド（下線なし）
% 分数用コマンド（表専用 - 分子を空欄にしてレイアウト保持）
\newcommand{\fracblank}[2]{%
  \ifshowanswer
    \textcolor{blue}{\frac{#1}{#2}}%  解答版：青色
  \else
    \frac{\phantom{#1}}{#2}%  空欄版：分子を透明化、レイアウト保持
  \fi
}

% enumerate環境のデフォルト設定（（１）、（２）形式）
\setlist[enumerate,1]{label=（\arabic*）, leftmargin=*, labelsep=0.5em}
\setlist[enumerate,2]{label=（\alph*）, leftmargin=*, labelsep=0.5em}

\usepackage[normalem]{ulem} % 下線用
\usepackage{stackengine} % 追加

% 丸囲み文字コマンド
\newcommand{\ctext}[1]{\raise0.2ex\hbox{\textcircled{\scriptsize{#1}}}}

% 回答行のラベル（伸ばし棒ではなく短い横線）
\newcommand{\answermark}{\rule[0.5ex]{1.1em}{0.5pt}}

% 定理スタイルのカスタマイズ（最後のピリオドを削除）
\newtheoremstyle{myplain}
  {3pt}   % Space above
  {3pt}   % Space below
  {\itshape}  % Body font
  {}  % Indent amount
  {\bfseries} % Theorem head font
  {}  % Punctuation after theorem head (空にしてピリオドを削除)
  {.5em}  % Space after theorem head
  {}  % Theorem head spec

\newtheoremstyle{mydefinition}
  {3pt}   % Space above
  {3pt}   % Space below
  {}  % Body font
  {}  % Indent amount
  {\bfseries} % Theorem head font
  {}  % Punctuation after theorem head (空にしてピリオドを削除)
  {.5em}  % Space after theorem head
  {}  % Theorem head spec

\theoremstyle{mydefinition}
\newtheorem{definition}{Definition}
\newtheorem*{definition*}{Definition}
\theoremstyle{myplain}
\newtheorem{theorem}{Theorem}
\newtheorem*{theorem*}{Theorem}

\newcounter{technique}

% 定理・定義番号を subsection ごとにリセット
\numberwithin{theorem}{subsection}
\numberwithin{definition}{subsection}
\numberwithin{technique}{subsection}

% 定理・定義番号を「1.1.1」形式に
\renewcommand{\thetheorem}{\thesubsection.\arabic{theorem}}
\renewcommand{\thedefinition}{\thesubsection.\arabic{definition}}
\renewcommand{\thetechnique}{\thesubsection.\arabic{technique}}

% 式番号を「1.1」形式に（sectionごとにリセット）
\numberwithin{equation}{section}
\renewcommand{\theequation}{\arabic{section}.\arabic{equation}}

% ========================================
% 問題環境の設定
% ========================================
% 問題番号カウンター（単純な連番: 1, 2, 3, ...）
\newcounter{problemcounter}
\renewcommand{\theproblemcounter}{\arabic{problemcounter}}

% 問題環境（枠なし、シンプル）
\newenvironment{problem}[1][\relax]{%
  \refstepcounter{problemcounter}%
  \vspace{2em}%
  \noindent%
  \textbf{問題\theproblemcounter}%
  \ifx#1\relax\else~\textbf{#1}\fi%
  \par\vspace{0.8em}%
  \noindent%
}{%
  \par\vspace{2em}%
}

% 問題環境（枠付き）
\newenvironment{problembox}[1][\relax]{%
  \refstepcounter{problemcounter}%
  \vspace{2em}%
  \begin{itembox}[l]{\textbf{問題\theproblemcounter}%
  \ifx#1\relax\else~\textbf{#1}\fi}%
}{%
  \end{itembox}%
  \vspace{1.5em}%
}

% 例題環境（枠付き）- 【】内の番号を使用
\newenvironment{example}[1][\relax]{%
  \vspace{2em}%
  \begin{itembox}[l]{\textbf{例題}%
  \ifx#1\relax\else~\textbf{#1}\fi}%
}{%
  \end{itembox}%
  \vspace{1.5em}%
}

% 練習問題環境（枠付き）- 【】内の番号を使用
\newenvironment{exercise}[1][\relax]{%
  \vspace{2em}%
  \begin{itembox}[l]{\textbf{練習問題}%
  \ifx#1\relax\else~\textbf{#1}\fi}%
}{%
  \end{itembox}%
  \vspace{1.5em}%
}

% 解答環境
\newenvironment{solution}{%
  \vspace{1em}%
  \noindent%
  \ifshowanswer%
    \textbf{【解答】}\\%
  \fi%
}{%
  \par\vspace{1em}%
}

% 問題専用の解答環境（常に空欄、手書きスペース確保）
% 使い方: \begin{problemsolution}[高さ] ... \end{problemsolution}
% デフォルト高さは15cm
% ※最後の空白は手動で \vspace{5\baselineskip} を追加してください
% 空欄時: 解答内容を白色で出力し、証明終了記号（□）を同じ位置に配置
\NewEnviron{problemsolution}[1][15cm]{%
  \begin{proof}\mbox{}\\%
  \ifshowanswer%
    {\color{blue}\BODY}%
  \else%
    {\color{white}\BODY}%
  \fi%
  \end{proof}%
}

% 別解環境（解答表示時のみ表示、ピリオドなし）
\NewEnviron{altproof}{%
  \ifshowanswer%
    \par\vspace{1em}%
    \noindent\textbf{【別解】}\par%
    {\color{blue}\BODY\par\hfill$\square$}%
    \vspace{1em}%
  \fi%
}

% 証明環境内の隠し解答ブロック
% 使い方: \begin{hiddenanswer} ... \end{hiddenanswer}
% 空欄時: 改ページしない範囲で \vfill により手書きスペースを確保 + 白文字でレイアウト保持
% 解答版: 青文字で表示（preamble の \showanswertrue で切り替え）
\NewEnviron{hiddenanswer}{%
  \ifshowanswer
    {\color{blue}\BODY}%
  \else
    \par\vfill
    {\color{white}\BODY}%
  \fi
}

% ========================================
% 穴埋めコマンド（\showanswer で表示切替）
% ========================================
% 使い方:
%   \fitblank[]{答え}        … 答え幅+計4文字分の下線（文中用）
%   \fitblank[答え]{}        … 同上（数学教材と同じ書き方）
%   \fitblank[答え]{幅}      … 幅（em）を指定
%   \fitblank{幅}            … 答えなしの空欄のみ
\makeatletter
% 文中穴埋め（\fitblank[]{答え}）の下線：答えの幅 + 両側2em（計4文字分の余白）
\newcommand{\fitblank@underlinewithpad}[1]{%
  \underline{\hspace{2em}#1\hspace{2em}}%
}
\newcommand{\fitblank@showinline}[1]{%
  \ifshowanswer
    \fitblank@underlinewithpad{{\color{blue}#1}}%
  \else
    \fitblank@underlinewithpad{\phantom{#1}}%
  \fi
}
\newcommand{\fitblank@showbfinline}[1]{%
  \ifshowanswer
    \fitblank@underlinewithpad{{\color{blue}\textbf{#1}}}%
  \else
    \fitblank@underlinewithpad{\phantom{\textbf{#1}}}%
  \fi
}
\newcommand{\fitblank@showwide}[1]{%
  \ifshowanswer
    \underline{\hspace{3em}{\color{blue}\ensuremath{\bm{#1}}}\hspace{3em}}%
  \else
    \underline{\hspace{3em}\phantom{\ensuremath{#1}}\hspace{3em}}%
  \fi
}
\newcommand{\fitblank@showbfwide}[1]{%
  \ifshowanswer
    \underline{\hspace{3em}{\color{blue}\ensuremath{\bm{#1}}}\hspace{3em}}%
  \else
    \underline{\hspace{3em}\phantom{\ensuremath{#1}}\hspace{3em}}%
  \fi
}
\newcommand{\fitblank@showwithwidth}[2]{%
  \ifshowanswer
    \underline{\hspace{3em}\makebox[#2em][c]{\color{blue}\ensuremath{\bm{#1}}}\hspace{3em}}%
  \else
    \underline{\hspace{3em}\makebox[#2em][c]{\phantom{\ensuremath{#1}}}\hspace{3em}}%
  \fi
}
\newcommand{\fitblank@resolve}[2]{%
  \def\fitblank@opt{#1}%
  \def\fitblank@mand{#2}%
  \ifx\fitblank@mand\empty
    \fitblank@showwide{#1}%
  \else
    \ifx\fitblank@opt\empty
      \fitblank@showinline{#2}%
    \else
      \fitblank@showwithwidth{#1}{#2}%
    \fi
  \fi
}
\newcommand{\fitblank@resolvebf}[2]{%
  \def\fitblank@opt{#1}%
  \def\fitblank@mand{#2}%
  \ifx\fitblank@mand\empty
    \fitblank@showbfwide{#1}%
  \else
    \ifx\fitblank@opt\empty
      \fitblank@showbfinline{#2}%
    \else
      \fitblank@showwithwidth{#1}{#2}%
    \fi
  \fi
}
\newcommand{\fitblank}[2][\relax]{%
  \ifx#1\relax
    \underline{\hspace*{#2em}}%
  \else
    \fitblank@resolve{#1}{#2}%
  \fi
}
\newcommand{\fitblankbf}[2][\relax]{%
  \ifx#1\relax
    \underline{\hspace*{#2em}}%
  \else
    \fitblank@resolvebf{#1}{#2}%
  \fi
}

% 会話例リスト用：下線の始点・長さを揃える（\fitblank[]{英語} を複数行に並べる場合）
\newlength{\dialogueblankindent}
\setlength{\dialogueblankindent}{2em}
\newcommand{\dialogueblankwidth}{32em}
\newcommand{\dialogueblank@print}[2]{%
  \ifshowanswer
    \underline{\makebox[#1][l]{\hspace{2em}{\color{blue}#2}\hspace{2em}}}%
  \else
    \underline{\makebox[#1][l]{\hspace{2em}\phantom{#2}\hspace{2em}}}%
  \fi
}
\newcommand{\dialogueblank@printcolumn}[2]{%
  \ifshowanswer
    \underline{\makebox[#1][l]{\hspace{1em}{\color{blue}\small #2}\hspace{1em}}}%
  \else
    \underline{\makebox[#1][l]{\hspace{1em}{\small\phantom{#2}}\hspace{1em}}}%
  \fi
}
\newcommand{\dialogueblank}[1]{%
  \par\vspace{0.3em}%
  \noindent\hspace{\dialogueblankindent}%
  \dialogueblank@print{\dialogueblankwidth}{#1}%
}
% 矢印の直後など、同じ行に続けて使う版
\newcommand{\dialogueblankinline}[1]{%
  \dialogueblank@print{\dialogueblankwidth}{#1}%
}
% minipage 等の半分カラム用（\linewidth に合わせた下線）
\newcommand{\dialogueblankcolumn}[1]{%
  \par\vspace{0.3em}%
  \noindent
  \dialogueblank@printcolumn{\linewidth}{#1}%
  \par\vspace{0.8em}%
}
\newcommand{\fitblankfixed}[2][\relax]{%
  \ifx#1\relax
    \underline{\hspace*{#2em}}%
  \else
    \fitblank@resolve{#1}{#2}%
  \fi
}
\newcommand{\fitblankblue}[2][\relax]{%
  \ifx#1\relax
    \underline{\hspace*{#2em}}%
  \else
    \def\fitblank@opt{#1}%
    \def\fitblank@mand{#2}%
    \ifx\fitblank@mand\empty
      \ifshowanswer
        \fitblank@underlinewithpad{{\color{blue}#1}}%
      \else
        \fitblank@underlinewithpad{\phantom{#1}}%
      \fi
    \else
      \ifx\fitblank@opt\empty
        \ifshowanswer
          \fitblank@underlinewithpad{{\color{blue}#2}}%
        \else
          \fitblank@underlinewithpad{\phantom{#2}}%
        \fi
      \else
        \ifshowanswer
          \textcolor{blue}{\underline{\makebox[#2em][c]{#1}}}%
        \else
          \underline{\hspace*{#2em}}%
        \fi
      \fi
    \fi
  \fi
}
\newcommand{\fitblankbfblue}[2][\relax]{%
  \ifx#1\relax
    \underline{\hspace*{#2em}}%
  \else
    \def\fitblank@opt{#1}%
    \def\fitblank@mand{#2}%
    \ifx\fitblank@mand\empty
      \ifshowanswer
        \fitblank@underlinewithpad{{\color{blue}\textbf{#1}}}%
      \else
        \fitblank@underlinewithpad{\phantom{\textbf{#1}}}%
      \fi
    \else
      \ifx\fitblank@opt\empty
        \ifshowanswer
          \fitblank@underlinewithpad{{\color{blue}\textbf{#2}}}%
        \else
          \fitblank@underlinewithpad{\phantom{\textbf{#2}}}%
        \fi
      \else
        \ifshowanswer
          \textcolor{blue}{\underline{\makebox[#2em][c]{\textbf{#1}}}}%
        \else
          \underline{\hspace*{#2em}}%
        \fi
      \fi
    \fi
  \fi
}
\makeatother

% 穴埋めコマンドの修正版（upLaTeX対応）
\newcommand{\Blank}[2][3cm]{%
  ~%
  \ifshowanswer
    \stackon[1pt]{\uline{\makebox[#1][l]{\hphantom{あ}}}}{{\normalsize\sffamily #2}}%
  \else
    \uline{\makebox[#1][l]{\hphantom{あ}}}%
  \fi
  ~%
}


% ========================================
% tcolorbox環境（背景パターンで区別）
% ========================================

% 定義環境（赤系：薄い赤背景 + 赤枠）- 自動番号付き
% 注: #1はラベル（オプション）、#2はタイトル（必須）
\newtcolorbox{definitionbox}[2][]{
    enhanced,
    colback=red!5,
    colframe=red!70!black,
    boxrule=0.8pt,
    sharp corners,
    left=10pt, right=10pt, top=8pt, bottom=8pt,
    fonttitle=\bfseries,
    code={\refstepcounter{definition}\if\relax\detokenize{#1}\relax\else\label{#1}\fi}, % #1が空でない場合のみラベル設定
    title={Definition \thedefinition\if\relax\detokenize{#2}\relax\else\quad\textbf{#2}\fi}
}

% 定理環境（青系：薄い青背景 + 青枠）- 自動番号付き
% 注: #1はラベル（オプション）、#2はタイトル（必須）
\newtcolorbox{theorembox}[2][]{
    enhanced,
    colback=blue!5,
    colframe=blue!70!black,
    boxrule=1pt,
    sharp corners,
    left=10pt, right=10pt, top=8pt, bottom=8pt,
    fonttitle=\bfseries,
    code={\refstepcounter{theorem}\if\relax\detokenize{#1}\relax\else\label{#1}\fi}, % #1が空でない場合のみラベル設定
    title={Theorem \thetheorem\if\relax\detokenize{#2}\relax\else\quad\textbf{#2}\fi}
}

% 問題環境（濃い灰色背景 + 太線）
\newtcolorbox{problembox2}[1][]{
    enhanced,
    colback=black!10,
    colframe=black,
    boxrule=1.2pt,
    sharp corners,
    left=10pt, right=10pt, top=8pt, bottom=8pt,
    fonttitle=\bfseries,
    title=#1
}

% テクニック環境（緑系：薄い緑背景 + 緑枠）- 自動番号付き
% 注: #1はラベル（オプション）、#2はタイトル（必須）
\newtcolorbox{techniquebox}[2][]{
    enhanced,
    colback=green!5,
    colframe=green!70!black,
    boxrule=0.8pt,
    sharp corners,
    left=10pt, right=10pt, top=8pt, bottom=8pt,
    fonttitle=\bfseries,
    code={\refstepcounter{technique}\if\relax\detokenize{#1}\relax\else\label{#1}\fi},
    title={Technique \thetechnique\if\relax\detokenize{#2}\relax\else\quad\textbf{#2}\fi}
}

% ========================================
% 手書き用空欄コマンド（下線なし、スペース確保）
% ========================================

% 解答ブロック全体の表示/非表示
% 使い方: \answerblock{解答内容}
\newcommand{\answerblock}[1]{%
  \ifshowanswer
    \textcolor{blue}{#1}%  解答を青色で表示
  \else
    % 空欄時は何も表示しない
  \fi
}

% ========================================
% answer用コマンド（レイアウト基準は答えの幅、空欄時も幅を確保）
% ========================================

% インライン用（テキスト・数式両用）
% 使い方: \answertext{答え}
% 空欄時: 答えの幅 + 両端に1em（1文字分）の余白
\newcommand{\answertext}[1]{%
  \ifshowanswer
    \textcolor{blue}{#1}%  解答を青色で表示
  \else
    \hspace{1em}\phantom{#1}\hspace{1em}%  空欄時は答えの幅 + 両端1em
  \fi
}

% 証明環境の右下用の最終回答（証明ブロック内・右端から 3em 内側）
% 使い方: \finalanswer{答え：220円}
\newcommand{\finalanswer}[1]{%
  \par\vspace{0.5em}%
  \noindent\hbox to \linewidth{%
    \hss
    \ifshowanswer
      \textcolor{blue}{#1}%
    \else
      \phantom{#1}%
    \fi
    \hspace{3em}%
  }%
  \par
}

% 数式ブロック用（align環境など内部で使用）
% 使い方: \answermath{数式}
% 空欄時: 答えの幅 + 両端に2em（2文字分）の余白
\newcommand{\answermath}[1]{%
  \ifshowanswer
    {\color{blue}#1}%  数式を青色で表示
  \else
    \hspace{2em}\phantom{#1}\hspace{2em}%  空欄時は答えの幅 + 両端2em
  \fi
}

% 表専用コマンド（空欄にしてレイアウト保持）
% 使い方: \answertable{答え}
% 空欄版：透明化してレイアウト保持
\newcommand{\answertable}[1]{%
  \ifshowanswer
    \textcolor{blue}{#1}%  解答版：青色で表示
  \else
    \phantom{#1}%  空欄版：透明化、レイアウト保持
  \fi
}

% ========================================
% blank用コマンド（解答時も余白あり版）
% ※通常は answer用コマンドを使用してください
% ========================================

% インライン用（テキスト・数式両用）
% 使い方: \blanktext{答え}
% 解答時・空欄時とも: 答えの幅 + 両端に0.5em空白
\newcommand{\blanktext}[1]{%
  \ifshowanswer
    \hspace{0.5em}\textcolor{blue}{#1}\hspace{0.5em}%
  \else
    \hspace{0.5em}\phantom{#1}\hspace{0.5em}%
  \fi
}

% 数式ブロック用（align環境など内部で使用）
% 使い方: \blankmath{数式}
% 解答時・空欄時とも: 答えの幅 + 左1.5em、右2.5em空白
\newcommand{\blankmath}[1]{%
  \ifshowanswer
    \hspace{1.5em}{\color{blue}#1}\hspace{2.5em}%
  \else
    \hspace{1.5em}\phantom{#1}\hspace{2.5em}%
  \fi
}

% 手書きスペース確保版（指定した高さの空白を作る）
% 使い方: \answerspace[高さ]{解答内容}
% 高さのデフォルトは5cm
\newcommand{\answerspace}[2][5cm]{%
  \ifshowanswer
    \textcolor{blue}{#2}%  解答を青色で表示
  \else
    \vspace{#1}%  指定した高さの空白を確保
  \fi
}

% fitblankblue / fitblankbfblue は preamble 前半で定義

% 分数の青色版
% 分数用コマンド（表専用 - 分子を空欄にしてレイアウト保持）
\newcommand{\fracblankblue}[2]{%
  \ifshowanswer
    \textcolor{blue}{\frac{#1}{#2}}%  解答版：青色
  \else
    \frac{\phantom{#1}}{#2}%  空欄版：分子を透明化、レイアウト保持
  \fi
}

